\label{cap:introducao_modelo}
A contextualização do problema de pesquisa e as premissas sobre os desafios no ensino e na aprendizagem de programação, que fundamentam esta introdução e suas subseções, baseiam-se nas investigações de \citeonline{souza2016}, \citeonline{stefik2013} e \citeonline{ferreira2025}.

A habilidade de programar computadores tornou-se uma competência essencial na sociedade contemporânea. Contudo, o ensino e a aprendizagem de algoritmos representam tarefas historicamente complexas, marcadas por altos índices de reprovação e evasão nos anos iniciais dos cursos de tecnologia. Um dos principais fatores que contribuem para esse cenário é a dificuldade dos alunos em compreender os conceitos abstratos da lógica computacional enquanto, simultaneamente, precisam lidar com a rigidez sintática das linguagens de programação tradicionais.

Em linguagens comerciais de mercado, um simples erro de digitação, a ausência de um ponto e vírgula ou o desconhecimento de uma palavra-chave reservada são suficientes para interromper a compilação, gerando mensagens de erro frequentemente incompreensíveis para um novato. Essa exigência de memorização gramatical atua como uma barreira inicial severa, desviando o foco do estudante, que deveria estar concentrado no desenvolvimento do raciocínio algorítmico, para a mera adequação de formatação textual.

Neste contexto, o presente trabalho parte do seguinte questionamento de pesquisa: de que forma a flexibilização e a personalização da sintaxe de uma linguagem de programação podem reduzir a barreira linguística e cognitiva no ensino introdutório de algoritmos?

Para responder a essa questão, aventa-se a seguinte hipótese: se o estudante iniciante tiver a liberdade de modelar os comandos e as regras estruturais da linguagem de acordo com o seu próprio idioma e curva de aprendizagem, haverá uma redução significativa da frustração associada a erros gramaticais, permitindo uma assimilação mais fluida e direta da lógica de programação.

\section{Justificativa}

A justificativa para a realização desta pesquisa reside na lacuna pedagógica e tecnológica observada nas atuais ferramentas de apoio educacional. Embora existam excelentes plataformas e linguagens adaptadas para iniciantes, a esmagadora maioria ainda impõe uma gramática estática e inalterável. 

Desenvolver uma ferramenta que inverta esse paradigma — adaptando a linguagem ao aluno, e não o aluno à linguagem — representa uma contribuição inovadora para a área de informática na educação. O interpretador dinâmico proposto ataca diretamente a carga cognitiva excessiva, configurando-se como um laboratório seguro de experimentação que facilita a posterior transição do estudante para as linguagens estritas exigidas pelo mercado de trabalho e por juízes online.

\section{Objetivos}

A seguir, são apresentados o objetivo geral e os objetivos específicos que nortearam o desenvolvimento desta pesquisa e a construção da ferramenta de software.

\section{Objetivo Geral}

O objetivo geral deste trabalho é projetar e desenvolver um interpretador dinâmico integrado a uma IDE web que permita a personalização da sintaxe e dos comandos de programação, visando atuar como uma ferramenta facilitadora no ensino e na aprendizagem de algoritmos para estudantes iniciantes.

\section{Objetivos Específicos}

Para viabilizar o alcance do objetivo geral, foram definidos os seguintes objetivos específicos e metodológicos: estudar os fundamentos de construção de compiladores e analisar ferramentas correlatas na literatura; projetar e implementar os módulos de análise léxica, sintática e semântica com suporte à injeção dinâmica de regras gramaticais; desenvolver uma interface web intuitiva para a interação do usuário com o interpretador; e, por fim, validar a plataforma através de estudos de caso para avaliar a sua viabilidade técnica e pedagógica.

\section{Estrutura do Trabalho}

Para propiciar ao leitor uma visão geral do que será apresentado, este documento está organizado da seguinte maneira: o Capítulo 2 estabelece o referencial teórico, abordando os conceitos de compiladores, metodologias de software e os trabalhos relacionados. O Capítulo 3 descreve a metodologia e os materiais, apresentando a modelagem da arquitetura da solução proposta. O Capítulo 4 detalha os estudos de caso e a implementação prática do interpretador. Por fim, o Capítulo 5 apresenta os resultados obtidos, culminando no Capítulo 6, que expõe as conclusões da pesquisa e sugere direcionamentos para trabalhos futuros.
